\documentclass[11pt,a4paper]{article}

\usepackage[utf8]{inputenc}

\usepackage[T1]{fontenc}

\usepackage[polish]{babel}

\usepackage{lmodern}

\usepackage{geometry}

\geometry{margin=27mm}

\usepackage[hidelinks]{hyperref}

\title{Triangulacja ocen rodziców i nauczycieli}

\author{Mateusz Pomianek}

\date{Szkic naukowy, wersja 1.0.0 | 28 listopada 2026}

\begin{document}

\maketitle

\begin{abstract}

Proponowana miara główna to różnica średnich bezwzględnych błędów predykcji.

\end{abstract}

\section{Wprowadzenie}

Raporty rodziców i nauczycieli dostarczają informacji z różnych kontekstów (De Los Reyes et al., 2015). Opis czasu i sytuacji pomaga interpretować zmienność pomiaru zachowania (Shiffman et al., 2008).

\section{Cel i pytania badawcze}

Hipotezy robocze wymagają zatwierdzenia przed rozpoczęciem właściwej weryfikacji.

\paragraph{H1.} Model wykorzystujący oba raporty poprawi predykcję niezależnej obserwacji zachowania względem lepszego z modeli jednoźródłowych w walidacji na nowych dzieciach.

\paragraph{H2.} Przyrost będzie większy dla kryterium obserwowanego równolegle w domu i w szkole niż dla kryterium ograniczonego do jednego środowiska.

\section{Materiał i metody}

Pierwszorzędowy wskaźnik: różnica średnich bezwzględnych błędów predykcji.

Ocenę dorosłych należy zebrać w tym samym tygodniu, a niezależną obserwację wykonać bez dostępu do raportów. Główna miara to błąd bezwzględny prognozy ujednoliconego wyniku obserwacji. Niekompletne pary i współzależność dzieci w klasach wymagają analizy wrażliwości.

Przed badaniem z udziałem ludzi należy zatwierdzić wymagania etyczne, zasady zgody i ochronę danych.

\section{Plan analizy}

Zamrożenie estymandu, reguł oceny jakości i obsługi braków nastąpi przed dostępem do danych wynikowych. Raport obejmie oszacowanie kontrastu głównego z przedziałem niepewności i pełne mianowniki. Jeśli przewidziano hipotezy drugorzędowe, przed analizą należy ustalić sposób kontroli wielokrotności.

\section{Wyniki}

W niniejszym pliku nie umieszczono wyników empirycznych, statystyk ani nieprzeprowadzonych analiz. Sekcja wymaga uzupełnienia po uzyskaniu i weryfikacji materiału dowodowego.

\section{Dyskusja i ograniczenia}

Ograniczenia obejmują możliwe błędy pomiaru, dobór próby, alternatywne mechanizmy i ograniczoną możliwość uogólniania. Korelacja nie dowodzi przyczynowości; poprawność techniczna nie jest dowodem znaczenia klinicznego.

\section{Wnioski}

Hipotezy i sposób pomiaru stanowią podstawę oceny wykonalności oraz przygotowania właściwego protokołu.

\section*{Bibliografia}

{\setlength{\parindent}{0pt}
\setlength{\parskip}{5pt}
\hangindent=1.27cm
De Los Reyes, A., Augenstein, T. M., Wang, M., Thomas, S. A., Drabick, D. A. G., Burgers, D. E., \& Rabinowitz, J. (2015). The validity of the multi-informant approach to assessing child and adolescent mental health. \textit{Psychological Bulletin, 141}(4), 858--900. \url{https://doi.org/10.1037/a0038498}\par
\hangindent=1.27cm
Shiffman, S., Stone, A. A., \& Hufford, M. R. (2008). Ecological momentary assessment. \textit{Annual Review of Clinical Psychology, 4}, 1--32. \url{https://doi.org/10.1146/annurev.clinpsy.3.022806.091415}\par
}

\end{document}
